\section{Centres, density, and prime heights}
\label{sec:centres}

Fix one of the finite blocks from Section~\ref{sec:preliminaries}, of
mass $w\le W\le 2w\le1$. All rows in this section belong to that block,
and their positive spatial widths are exact dyadic numbers. We assign
rows of one width to integer multiples of their denominators. The
assignment will have two complementary properties: an assigned-row
density bound at small spatial scale, and a quantitative separation between families
assigned at widely different scales. These properties supply the
summable centre kernels used in the later pair counts.

\subsection{Popular multiples and the density bound}

Fix $H_0=.001$. A small constant $k_0>0$ will be chosen after the other
fixed parameters. Write
\[
 W_r=\sum_{\substack{v\text{ in the block}\\r(v)=r}}\phi(v)r.
\]

\begin{definition}[Popular multiples and original hosts]
\label{cen:popular}
An integer $D\ge1$ is \emph{popular at width $r$} if $Dr\le k_0$ and
\[
 \sum_{\substack{v\mid D\\r(v)=r}}\frac{\phi(v)}D
 \ge (Dr)^{H_0}.
\]
A row $v$ is \emph{raw} if it divides no popular multiple at its width.
Otherwise its \emph{original host} is the largest popular multiple at
that width divisible by $v$.
\end{definition}

The largest host exists: at a fixed positive width, all popular
multiples are at most $k_0/r$.

\begin{lemma}[Density above original hosts]
\label{cen:density}
Let $\mathcal V$ be a subfamily of rows of a common width $r$, all
dividing an integer $L$. Suppose that $L$ is at least the original host
of every nonraw row in $\mathcal V$; raw rows are allowed as well. Define
\[
 k_L=Lr,\qquad
 \delta_L=\sum_{v\in\mathcal V}\frac{\phi(v)}L,\qquad
 m_L=k_L\delta_L.
\]
Then
\begin{equation}
 \delta_L\ll_{k_0} k_L^{H_0}.
 \label{eq:centre-density}
\end{equation}
Moreover $\delta_L\le1$.
\end{lemma}

\begin{proof}
The last assertion follows from $\sum_{v\mid L}\phi(v)=L$.
If the family is nonempty, $2L$ cannot be popular at width $r$: any
assigned nonraw row would then have an original host at least $2L$,
whereas any assigned raw row would no longer be raw. If $2Lr\le k_0$,
failure of popularity gives
\[
 \frac12\delta_L
 \le\sum_{\substack{v\mid2L\\r(v)=r}}\frac{\phi(v)}{2L}
 <(2Lr)^{H_0}.
\]
If $2Lr>k_0$, use $\delta_L\le1$. The empty family is immediate.
\end{proof}

A \emph{centre family} consists of an integer $L$, a common width $r$,
and an assigned subfamily satisfying the hypotheses of
Lemma~\ref{cen:density}. For two such families, indexed by $i,j$, use the notation
\[
 L=L_i,\quad M=L_j,\quad G=(L,M),\quad U=L/G,\quad T=M/G,
\]
\[
 \begin{gathered}
 r_{\min}=\min(r_i,r_j),\qquad r_{\max}=\max(r_i,r_j),\\
 \Theta_{ij}=G r_{\min},\qquad K=[L,M]r_{\max}.
 \end{gathered}
\]
In particular $(U,T)=1$. The identity
\begin{equation}
 \Theta_{ij}
 =(m_Lm_M)^{1/2}(UT)^{-1/2}
 \left(\frac{r_{\min}}{r_{\max}}\right)^{1/2}
 (\delta_L\delta_M)^{-1/2}
 \label{eq:centre-kernel}
\end{equation}
will convert arithmetic separation into a summable mass kernel. We
apply it only to nonempty families, so its inverse densities are defined.

\begin{lemma}[Centre-kernel summation]
\label{cen:kernel}
Consider centre families satisfying Lemma~\ref{cen:density}, with at
most one family for each pair $(L,r)$ and total mass at most $W\le1$.
For every fixed $\alpha>1/2$ and sufficiently small fixed
$\epsilon>0$, there is $c>0$ such that
\[
 \sum_{i,j}\Theta_{ij}(\delta_{L_i}\delta_{L_j})^\alpha
 (U_{ij}T_{ij})^\epsilon
 \ll_{k_0,\alpha,\epsilon} W^{1+c}.
\]
\end{lemma}

\begin{proof}
Choose $u>1/2$ sufficiently close to $1/2$ that
$H_0(\alpha-u)>u-1/2$, and then $\epsilon<u-1/2$.
For a single centre,
\[
 m_L^{1/2-u}\delta_L^{\alpha-1/2}
 =k_L^{1/2-u}\delta_L^{\alpha-u}\ll_{k_0}1.
\]
Indeed use \eqref{eq:centre-density} when $k_L\le1$, and use
$\delta_L\le1$ when $k_L>1$. Hence the summand is at most
\[
 C_{k_0}(m_Lm_M)^u(UT)^{-1/2+\epsilon}
 \left(\frac{r_{\min}}{r_{\max}}\right)^{1/2}.
\]
At each pair of exact widths Lemma~\ref{ar:tensor} applies, since
$1/2-\epsilon>1-u$. If $M_r$ is the total mass at width $r$, the
remaining sum is bounded by
\[
 C\sum_{r,s}(M_rM_s)^u
 \left(\frac{\min(r,s)}{\max(r,s)}\right)^{1/2}
 \ll\sum_r M_r^{2u}\le W^{2u}.
\]
The middle inequality is convolution with a summable geometric kernel
on dyadic width indices. Take $c=2u-1$.
\end{proof}

\begin{lemma}[Raw rows suffice when they carry mass]
\label{cen:raw}
If raw rows have mass at least $w/4$ on arbitrarily late blocks, the
positive-measure avoided set fixed in Section~\ref{sec:preliminaries}
cannot exist.
\end{lemma}

\begin{proof}
For a raw row use centre $L=v$ and give every point of its grid weight
$R(v)=\phi(v)/v$. Then $\delta_L=R(v)$ and $m_L=\phi(v)r(v)$.
The determinant-cycle count of Lemma~\ref{ar:cycle} at distance
$O(r_{\max})$, multiplied by
$r_{\min}$ and the two row weights, is
\[
 O\bigl(m_Lm_M+\Theta_{ij}\delta_L\delta_M\bigr).
\]
Lemma~\ref{cen:density} applies to these raw singleton families.
Lemma~\ref{cen:kernel} with $\alpha=1$ therefore gives a uniform
second-moment bound for their weighted interval trains. Their means
are bounded below by a constant times $w$, and their integrals on each
fixed interval approach that interval's length times the mean: these
are uniformly weighted grids with $v\to\infty$ down the block tail.
Approximation of the avoided set by a finite union of intervals, using
the second-moment bound for the approximation error, gives a positive
integral on that set, a contradiction.
\end{proof}

We may consequently pass to arbitrarily late blocks on which the
nonraw rows carry a fixed positive fraction of the block mass. The
following adjustment of their hosts is made before any numerator
weights are imposed.

\subsection{Adjusted hosts and the density purge}

The original hosts give the density bound, but they do not yet separate
families assigned at different scales. We therefore allow a row to use a
popular centre from another width, charging for the required arithmetic
enlargement and width change. Minimizing this cost will ensure that a
centre at a different scale cannot offer a substantially cheaper fit;
Lemma~\ref{cen:exterior} makes this separation precise.

For a nonraw row $v$ of width $r$, minimize over all original popular
centres $(D,r_0)$ in the block the cost
\begin{equation}
 \begin{gathered}
 b+10a+10c_++(-c)_+,\\
 b=-\log(Dr_0),\quad
 a=\log\frac{v}{(v,D)},\quad c=\log\frac r{r_0},
 \end{gathered}
 \label{cen:assignment-cost}
\end{equation}
where $t_+=\max(t,0)$. Choose a minimizer, resolving ties
arbitrarily, and put $L=[D,v]=De^a$. Consolidate assignments with the
same triple
\[
 (L,r,B^*),\qquad B^*=\lfloor\text{minimum cost}\rfloor,
\]
into one \emph{centre object}; call the integer $B^*$ its \emph{shell}.
Write $i$ for such an object, and let
$B_i$ be the dyadic power of two with $B_i\le B_i^*<2B_i$.
We call $B_i$ its band. All costs, and hence all these bands, are large
when $k_0$ is small.

For each object define $k_{L_i},\delta_{L_i},m_{L_i}$ as in
Lemma~\ref{cen:density}, using precisely its assigned rows, and put
\[
 \mathcal F_i=\{L_i/v:v\text{ is assigned to }i\}.
\]
Thus $\pi_{L_i}(\mathcal F_i)=\delta_{L_i}$. Different objects
partition the assigned rows, including their masses. These symbols
will always refer to the original assignments, even when some masks
or numerator points are subsequently removed.

\begin{lemma}[Scale and number of adjusted centres]
\label{cen:adjusted}
Every adjusted centre satisfies Lemma~\ref{cen:density} and
\[
 |\log k_{L_i}|\le B_i^*+1,\qquad
 \log k_{L_i}\le .1(B_i^*+1).
\]
The number of original popular centres at a fixed width $r$ and with
$-\log(Dr)\in[b,b+1)$ is
\[
 \ll W_r e^{(1+3H_0)b}.
\]
The total number of adjusted centre objects with $B_i^*\le t$,
including all widths, is $\ll e^{2t}$. Deleting objects with
$\delta_{L_i}<e^{-.32B_i^*}$ loses $o(W)$ as $k_0\to0$.
\end{lemma}

\begin{proof}
For the chosen candidate in \eqref{cen:assignment-cost},
$\log k_L=-b+a+c$. Since $b>0$ and $a\ge0$, the cost bounds
$|\log k_L|$, and it is at least $10\log k_L$ when the latter is
positive. The cost is no larger than the cost obtained by using the
row's original host at its own width, namely minus the log of that
host's spatial scale. Since the cost also bounds $-\log k_L$, the
adjusted $L$ is at least that original host. This proves the first
assertions.

Let $n$ be the number of original popular centres in the indicated
unit $b$-range. Sum the popularity inequality after multiplying by
$Dr$, and apply Cauchy--Schwarz to row incidences. Using
$\sum_{v\mid g}\phi(v)=g$ gives
\[
 \bigl(n e^{-(1+H_0)b}\bigr)^2
 \ll W_r\sum_{D,D'}(D,D')r.
\]
Here $D,D'\asymp e^{-b}/r$, so Lemma~\ref{ar:tensor} gives, for $u>1/2$
arbitrarily close to $1/2$,
\[
 \sum_{D,D'}(D,D')r\ll e^{-b}n^{2u}.
\]
Consequently
\[
 n\ll W_r^{1/(2-2u)}
 \exp\left(\frac{1+2H_0}{2-2u}b\right)
 \ll W_r e^{(1+3H_0)b},
\]
after choosing $u$ sufficiently close to $1/2$ and using $W_r\le1$.

At a fixed source width $r_0$, sort the assignments into unit boxes
of $(b,a,c)$. For each source $D$, there are $O(e^a)$ possibilities
for the integer $v/(v,D)$ and hence for $L$. A unit $c$-range contains
only boundedly many dyadic widths, and a unit $(b,a,c)$-box gives
boundedly many integer shells. Thus the number of resulting objects
is
\[
 \ll W_{r_0}\exp((1+3H_0)b+a).
\]
Multiplication by $k_L=e^{-b+a+c}$ leaves the exponent
\[
 3H_0b+2a+c\le .21B_i^*+O(1).
\]
An object in the stated deletion contributes an additional factor
$e^{-.32B_i^*}$. This bounds its entire mass, not just the mass of
rows using the particular source witness. If an object has several
source witnesses, cover it by all their boxes, or choose any one:
the upper count can overcount objects but cannot omit one. The
coefficient inequality above holds also for $c<0$, since then
$c\le .21(-c)_+$. At shell $t$, the nonnegative cost bounds
$b,a,|c|$ by $O(t+1)$, so there are $O((t+1)^3)$ unit boxes.
All costs are at least $-\log k_0$. The deleted mass is therefore
at most
\[
 C\sum_{r_0}W_{r_0}
 \sum_{t\ge-\log k_0-O(1)}(t+1)^3e^{-.11t}=o(W).
\]
This also shows uniformity when a source-width mass is arbitrarily
small. Similarly, summing the object-count
bound over cost at most $t$, and using $\sum_{r_0}W_{r_0}\le1$,
gives $O(e^{2t})$.
\end{proof}

Retain the remaining objects. In addition to their upper density
bound, their original assigned families satisfy the following lower
bounds, before any subsequent mask or numerator removals:
\begin{equation}
 \delta_L\ge e^{-.32B_i^*},\qquad
 m_L\gg e^{-1.32B_i^*}.
 \label{eq:centre-lower}
\end{equation}

\subsection{Exterior separation and summation with height decay}

\begin{lemma}[Separation of different bands]
\label{cen:exterior}
If $B_i/B_j$ is sufficiently small, then every row $v$ assigned to
$j$ satisfies at least one of
\begin{equation}
 \frac{v}{(v,L_i)}\frac{r_j}{r_{\min}}>e^{.015B_j^*},
 \qquad
 \log(r_i/r_j)>.7B_j^*.
 \label{eq:exterior-gain}
\end{equation}
\end{lemma}

\begin{proof}
Choose a source $(D_i,r_{0i})$ that produced an assignment in object
$i$. Then $D_i\mid L_i$, and its source parameters in
\eqref{cen:assignment-cost} are $O(B_i)$. In particular
\[
 \log\frac{v}{(v,D_i)}
 \le \log\frac{v}{(v,L_i)}+\log\frac{L_i}{D_i}.
\]
Using $(D_i,r_{0i})$ as a candidate for $v$, and comparing the two
widths through $r_i$, bounds its minimum cost by
\[
 O(B_i)+10\log\frac{v}{(v,L_i)}
 +10\log\frac{r_j}{r_{\min}}+\log^+\frac{r_i}{r_j}.
\]
If both conclusions failed, this would be at most
$O(B_i)+.85B_j^*<B_j^*$, contrary to the definition of the shell.
\end{proof}

We record two ways in which \eqref{eq:centre-kernel} will absorb
errors. This also explains why the shell label in an adjusted centre
does not invalidate the earlier tensor summations.

\begin{lemma}[Width and height gains]
\label{cen:decay-kernels}
Under \eqref{eq:centre-lower}, the following kernels are summable
with a positive power of total mass beyond linear.
First, restrict to pairs with sufficiently small $B_i/B_j$ and
$\log(r_i/r_j)>.65B_j^*$. For sufficiently small absolute
$c_s,\epsilon>0$, the kernel is
\[
 \Theta_{ij}e^{c_sB_j}(UT)^\epsilon.
\]
Second, fix $c>0$ and $C>0$. Sum over dyadic $Y$ and pairs with
$B_i,B_j\le Y/A_0$. For $A_0$ sufficiently large the kernel is
\[
 \Theta_{ij}e^{-cY}\tau(UT)^C.
\]
The bounds permit the shell multiplicities of the adjusted objects
and the displayed height summation.
\end{lemma}

\begin{proof}
In \eqref{eq:centre-kernel}, changing the mass exponent from $1/2$
to $u>1/2$, including the inverse densities there, costs at most
\[
 C\exp\left(
  \bigl(.16+1.32(u-1/2)\bigr)(B_i^*+B_j^*)\right).
\]
For the first kernel, take $0<u-1/2\le .005$ and use band
separation so that $B_i^*/B_j^*\le .01$. The displayed conversion
cost is then at most $Ce^{.18B_j^*}$. Reserve the width power
$1/8$ for summation. Its remaining power $3/8$, together with
$\log(r_i/r_j)>.65B_j^*$, supplies
$e^{-.24375B_j^*}$. For $c_s\le .01$ this leaves at least
$e^{-.05B_j^*}$. Choose $\epsilon<u-1/2$, so the remaining
$UT$-price still meets the tensor inequality.

At a fixed pair of integer shells $s,t$, there is at most one
object at a given $(L,r)$. Tensor summation followed by convolution
in the widths, now with their power $1/8$, bounds this pair of
shells by $CW_s^uW_t^u\le CW^{2u}$, where $W_s,W_t$ are their
total masses. Thus the first kernel is bounded by
\[
 CW^{2u}\sum_{t\ge1}(t+1)e^{-.05t}\ll W^{2u}.
\]
There is no summation of an undamped uniform shell bound.

For the second kernel take $u=.75$, and write
$C_u=.16+1.32(u-1/2)$. Its shell-conversion cost is at most
$C\exp(4C_uY/A_0)$, since $B_i^*,B_j^*<2Y/A_0$.
Enlarge $A_0$ so that $4C_u/A_0<c/2$. The divisor bound
$\tau(UT)^C\ll_\epsilon(UT)^\epsilon$ uses less than the
available margin $u-1/2$ in the tensor inequality. At each pair
of shells, tensor and width summation again cost at most $CW^{2u}$.
There are $O((Y+1)^2)$ eligible shell pairs, so the total is at most
\[
 CW^{2u}\sum_{Y\text{ dyadic}}(Y+1)^2e^{-cY/2}
 \ll W^{2u}.
\]
Both exponents $2u$ are strictly greater than one.
\end{proof}

\subsection{Prime heights and starting scales}
\label{cen:height-starts}

Use dyadic heights $Y=1,2,4,\ldots$ for the prime bins
\[
 \mathcal P(Y)=\{p:Y\le\log\log p<2Y\}.
\]
Primes below these bins are always part of the low initial data.
A bin is \emph{dense on $L_i$} if
\[
 \sum_{\substack{p\mid L_i\\p\in\mathcal P(Y)}}\frac1p
 \ge\kappa Y.
\]
Let $Y_d(i)$ be the last dense height, or zero if there is none.
The two types and their starts are defined by
\[
 \begin{array}{c|c|c}
 \text{type}&\text{condition}&Y_0(i)\\ \hline
 \text{high}&Y_d(i)>D_0B_i&2Y_d(i)\\
 \text{medium}&Y_d(i)\le D_0B_i&\max(A_0B_i,2Y_d(i)).
 \end{array}
\]
All bins from the start onward are sparse. Colour the integers
$\log_2B_i$ modulo a fixed large integer. Retain a single type and
colour carrying the greatest mass. This keeps $\gg W$ mass, with a
fixed implied constant, and makes distinct retained bands as widely
separated as required below.

\subsection{Order of the fixed parameters}
\label{cen:hierarchy}

We give the order here because all subsequent selections must be
uniform down the block tail. Some constants name thresholds that
are used only later: their quantitative roles are specified below,
and their constructions are given at the indicated sections. Every
choice is fixed independently of the block. This is the dependency
order of the constants, not the chronology of the adaptive construction;
the list also serves as a reference for the later estimates.

\begin{enumerate}
\item \emph{Prime powers and geometry.}
Choose first a sufficiently large absolute $C_{\rm mask}$, bounding
the excess prime powers in a deficit mask. Choose a sufficiently
large absolute $\beta$, which will join centres at height $Y$ when
$\log UT\le\beta Y$. It must absorb the absolute prime-product
costs $e^{CY}$ in the distant comparisons of
Sections~\ref{sec:partitions} and~\ref{sec:joins}.
Set $\sigma_g=.1$. Choose a small dyadic $\sigma>0$ so that absolute
harmonic costs below height $\sigma Y$ are at most $e^{Y/100}$.
Choose
\[
 0<h_*\ll\sigma h_H,\qquad 0<h_H\ll\sigma,
 \qquad 3h_*<h_a\ll\sigma_g,
\]
with arbitrarily large separations. Here $h_*$ controls the logarithmic
sizes of comparison data within a component, $h_H$ controls the
rational-model heights, and $h_a$ bounds the number of offset grids
used to record very close pairs. Finally in this stage take the
dyadic $A_0$ sufficiently large. In particular, path costs
$e^{CY/A_0}$ must be at most $e^{h_*Y/10}$, and the height-decay
applications of Lemma~\ref{cen:decay-kernels} must hold.

\item \emph{Occupancy and band separation.}
Choose $s'>0$ sufficiently small absolutely for the later powers
of divisor counts used with \eqref{eq:exterior-gain}, then take
$s_{\rm low}>0$ sufficiently smaller for the mask-moment deletion.
Choose dyadic $\ell>0$ and $\kappa>0$ such that $\ell<\sigma$ and
both $\ell$ and $\kappa A_0$ are arbitrarily small relative to
$s_{\rm low}$. Next take dyadic $D_0$ large compared with
$1/\kappa$ and $A_0$. The colour separation is now taken large
compared with $D_0/\ell$, $D_0/(\sigma A_0)$, and the requirements
of \eqref{eq:exterior-gain} and the centre counts.

\item \emph{Rational models and localization.}
The direct rotation-switch budget $D'$ in Section~\ref{sec:direct}
is large compared with $\beta,D_0$ and the costs depending on
them. Choose the large model exponent $E$ after these resulting
linear exponential bounds, then $Q_0\gg E$ for the deep direct
comparisons, and then a still larger $J>4E$ for the tag cutoff
$e^{JB_i}$. Choose $E_g$ sufficiently large after these constants
for the birth spatial scale $r_i e^{-E_gB_i}$, and then choose
$E_2$ much larger than $E_g$ for the raw partition spacing.

For the low rational models and their soft tests, fix
$0<\ell_{\rm sh}<1/25$, then $0<h_0\ll\ell_{\rm sh}$,
$0<h_s<.01$, and finally $\epsilon_s>0$ sufficiently small relative
to $h_s$ and the measure $\mu$ of the avoided set.

The remaining choices in this stage control inverse estimates and
rotation switches. Small powers in
the inverse and period estimates can be chosen arbitrarily small
after the fixed comparisons. Choose the harmonic-mask cutoff
$d_0>0$ after the corresponding period cutoffs. In particular the
powers used for period costs and $\tau(UT)$ in mixed comparisons
are small relative to $\beta A_0$. Coefficients of harmonic
exponents to be absorbed by $s_{\rm low}$, $\kappa A_0$, or
$\beta$ are absolute; they do not arise from these tiny-power
constants.

The sparse-bin switch budget $D''$ in Section~\ref{sec:sparse}
need only be sufficiently large absolutely; a factor
$\tau(UT)^C$ is carried separately there. For the last direct
subcases choose a small signature exponent $h'>0$, a small switch
threshold $t_1$, a large fixed $J_1$, and a small rotation-budget
exponent $\epsilon'$ in their required order. Concretely these
choices permit $h'\ll h_0$, $t_1\ll h'/Q_0$,
$J_1\gg J/\ell_{\rm sh}$, and $C_1\epsilon'\ll t_1$, with
$C_1$ the absolute constant in Lemma~\ref{ar:rotation}. All are fixed before taking the
arbitrarily small saturation exponent $\vartheta$ and the
associated parameters of \eqref{eq:small-saving}.

\item \emph{Final cutoffs.}
Take $k_0$ sufficiently small after all these choices. Every band
and every bin at a fixed positive linear scale in a band is then
as large as needed; in particular $\ell B_i$ and $\sigma Y$ are
bin heights. A sufficiently large fixed $\lambda_0$ for the
small-scale arithmetic comparisons and the hole deletion may
still depend on $k_0$. Then choose $w$ sufficiently small after
the constants depending on these data. Finally choose a fixed
cap $N_0$ as large as necessary. The comparison bounds used for
clipping in Section~\ref{sec:tables} will be independent of this
cap.
\end{enumerate}

The objects retained so far have a common type and colour, disjoint
row assignments, the upper density bound \eqref{eq:centre-density},
and the lower bounds \eqref{eq:centre-lower}. We next specify their
initial numerator weights, without using any later adaptive table.
